\subsection{Bandwidth-Related Effects}
\label{sec:correlated_effects}
\label{sec:refraction}
\label{sec:bandwidth_mosaic_cap}

The anomalous intensity along \(2\theta\) for the \(m=0\), \(L=3\) feature in Fig.~\ref{fig:hk0_low_l_star} can be explained by the bandwidth forming an Ewald shell on a much smaller Bragg sphere. In this
regime the measured intensity is not set by an external Ewald-sphere
intersection alone. 

For each sampled wavelength
\(\lambda_j\), with spectral weight \(w_j\), the model evaluates
\begin{equation}
  k_0(\lambda_j)=\frac{2\pi}{\lambda_j},
  \qquad
  \tilde n(\lambda_j)
  =
  1-\delta_n(\lambda_j)+i\beta_n(\lambda_j).
  \label{eq:n_complex_lambda_main}
\end{equation}
Here \(\delta_n\) is the refractive decrement that shifts the internal phase
velocity, and \(\beta_n\) is the absorptive correction that gives the finite
penetration depth. The real decrement may be estimated from the electron
density \(\rho_e\) and classical electron radius \(r_e\), while the imaginary
term is the wavelength-dependent optical input used directly by the complex
wavevector calculation:
\begin{equation}
  \delta_n(\lambda_j)
  =
  \frac{r_e\lambda_j^2\rho_e}{2\pi},
  \qquad
  \beta_n(\lambda_j)
  =
  \operatorname{Im}\{\tilde n(\lambda_j)\}.
  \label{eq:delta_beta_lambda_main}
\end{equation}
Thus bandwidth changes both the Ewald-sphere radius and the optical constants
that determine refraction and attenuation.
For an incident wavevector in air, the grazing angle and azimuth are
\begin{equation}
  \sin\theta_i
  =
  \frac{k_{i,z}}{k_0(\lambda_j)},
  \qquad
  \tan\varphi_i
  =
  \frac{k_{i,x}}{k_{i,y}}.
  \label{eq:incident_angles_lambda_main}
\end{equation}

The transmitted angle and internal wavevector scale are
\begin{align}
  \theta_t(\lambda_j)
  &=
  \cos^{-1}\!\left[
  \frac{\cos\theta_i}{\operatorname{Re}\{\tilde n(\lambda_j)\}}
  \right],\\
  k'(\lambda_j)
  &=
  k_0(\lambda_j)
  \left[\operatorname{Re}\{\tilde n(\lambda_j)^2\}\right]^{1/2}.
  \label{eq:theta_kprime_lambda_main}
\end{align}
The in-plane internal components are fixed by \(\theta_t\) and \(\varphi_i\).
The remaining surface-normal component is complex. Using the closed-form
solution of the \(a+i b\) decomposition, define
\begin{align}
  a_t(\lambda_j)
  &=
  k'(\lambda_j)^2\sin^2\theta_t(\lambda_j),\\
  b(\lambda_j)
  &=
  k_0(\lambda_j)^2
  \operatorname{Im}\{\tilde n(\lambda_j)^2\}.
  \label{eq:at_b_lambda_main}
\end{align}
Then
\begin{align}
  \operatorname{Re}\{k_{t,z}(\lambda_j)\}
  &=
  \left[
  \frac{\sqrt{a_t(\lambda_j)^2+b(\lambda_j)^2}+a_t(\lambda_j)}{2}
  \right]^{1/2},\\
  \operatorname{Im}\{k_{t,z}(\lambda_j)\}
  &=
  \left[
  \frac{\sqrt{a_t(\lambda_j)^2+b(\lambda_j)^2}-a_t(\lambda_j)}{2}
  \right]^{1/2}.
  \label{eq:ktz_solution_lambda_main}
\end{align}
The branch is chosen so that \(\operatorname{Im}\{k_{t,z}\}\ge 0\), giving a
decaying wave inside the film. The same result is applied to the outgoing
internal wavevector using the scattered internal angle. Optical reversibility
then refracts the scattered ray back into air. The algebraic derivation and the
interface transmission factors are given in the Supporting Information.

The imaginary normal components set the propagation loss. For each accepted
candidate, define
\begin{equation}
  \kappa_i(\lambda_j)
  =
  \operatorname{Im}\{k_{t,z}(\lambda_j)\},
  \qquad
  \kappa_f(\lambda_j)
  =
  \operatorname{Im}\{k'_{t,z}(\lambda_j)\} .
  \label{eq:kappa_if_lambda_main}
\end{equation}
If \(L_{\mathrm{in}}\) and \(L_{\mathrm{out}}\) are the entry and exit
propagation lengths used for that candidate, the absorption/evanescence weight
is
\begin{equation}
  A_{\mathrm{prop}}(\lambda_j)
  =
  \exp\!\left[
  -2\kappa_i(\lambda_j)L_{\mathrm{in}}(\lambda_j)
  -2\kappa_f(\lambda_j)L_{\mathrm{out}}(\lambda_j)
  \right].
  \label{eq:imkz_pathlength_weight_main}
\end{equation}
This is the form used by the simulation: the path lengths set how far the wave
propagates in the film, and the imaginary parts of the internal normal
wavevectors set the damping per unit length. For a uniformly scattering film of thickness \(t\),
this may be written as a depth average,
\begin{equation}
  \bar A_{\mathrm{prop}}(\lambda_j)
  =
  \frac{1}{t}\int_0^t
  \exp\!\left[-2\left(\kappa_i(\lambda_j)+\kappa_f(\lambda_j)\right)\zeta\right]
  d\zeta
  =
  \frac{
  1-
  \exp\!\left[-2\left(\kappa_i(\lambda_j)+\kappa_f(\lambda_j)\right)t\right]}
  {2\left(\kappa_i(\lambda_j)+\kappa_f(\lambda_j)\right)t}.
  \label{eq:absorption_lambda_main}
\end{equation}
This attenuation is multiplied by the entrance and exit transmission factors and
the structure-factor intensity.

The detector image is therefore an incoherent sum over wavelength samples,
\begin{equation}
  I_{\mathrm{det}}(x,y)
  =
  \sum_j
  w_j\,
  I_{\mathrm{det}}
  \!\left[
  x,y;
  \lambda_j,
  \tilde n(\lambda_j)
  \right].
  \label{eq:detector_sum_lambda_main}
\end{equation}
The bandwidth has two roles: it thickens the reciprocal-space intersection
through \(k_0(\lambda_j)\), and it changes the refracted wave through
\(\tilde n(\lambda_j)\). 

For the \(00L\) family, \(G_{00L}=2\pi L/c\). A fixed wavelength spread gives
an effective Ewald-shell thickness, and its angular overlap with the
mosaic-populated Bragg cap scales approximately as \(1/L\). Low-\(L\) Bragg
spheres therefore retain overlap with the mosaic cap over a broader range of
incident-angle detuning than high-\(L\) Bragg spheres. The long mosaic tail
supplies the tilted crystallites, while wavelength-dependent refraction and
attenuation determine how much of that intensity reaches the detector.

\begin{figure}[H]
  \centering

  \begin{minipage}[t]{0.48\linewidth}
    \centering
    \draftfigureplaceholder{figures/mosaic/special_cause_reciprocal_matrix_ratio_scaled.png}{%
      Bragg-sphere size comparison for \(L=3,6,9\), showing how a fixed
      wavelength bandwidth maps to a larger angular band for smaller \(L\).
    }
  \end{minipage}
  \hfill
  \begin{minipage}[t]{0.48\linewidth}
    \centering
    \draftfigureplaceholder{figures/mosaic/with_ewald.png}{%
      Ewald sphere with finite bandwidth thickness intersecting the mosaic cap
      of a low-\(L\) \(00L\) Bragg sphere.
    }
  \end{minipage}

  \caption[Bandwidth-related effects for the low-L feature]{Bandwidth-related effects for the low-\(L\)
  \(m=0\) feature. The wavelength distribution changes both the Ewald-sphere
  radius and the refractive index used to calculate the internal wavevector.
  Because \(G_{00L}\propto L\), a fixed bandwidth samples a larger angular
  fraction of low-\(L\) Bragg spheres than high-\(L\) Bragg spheres.}
  \label{fig:correlated_bragg_sphere_size_series}
\end{figure}

\subsection{Reflectivity}
\label{sec:reflectivity}

The bright near-origin \(m=0\) intensity in Fig.~\ref{fig:hk0_low_l_star} lies in the near-critical specular region.

The specular-family profiles in this near-critical region are calculated with an
optical-to-kinematic handoff. The low-\(Q_z\) branch is calculated with a Parratt
multilayer reflectivity recursion, which retains refraction, evanescence, and
multiple reflection at the film interfaces.\cite{Parratt1954} At larger
\(Q_z/Q_c\), where critical-angle corrections become small, this optical branch
is matched to a normalized finite-stack kinematic term. In the calculation, the
two limits are joined with a smooth log-intensity blend over the region where
they agree. Figure~\ref{fig:parratt_handoff} overlays the measured near-critical
\(m=0\) profile with the Parratt and kinematic limiting calculations on the same
axes, while omitting the internal blended curve. The direct overlay shows which
limit describes each part of the measured profile and where the two calculated
branches converge. This construction prevents near-critical specular intensity
from being treated as ordinary high-\(Q_z\) kinematic scattering while
preserving the finite-stack modulation that controls the higher-\(L\) rod
profile. Details of the recursion, scale estimate, and blend criterion are
provided in the Supporting Information.

\begin{figure}[H]
  \centering
  \includegraphics[width=0.98\linewidth]{figures/correlated_effects/fig_reflectivity_data_parratt_kinematic.png}
  \caption[Measured profile overlaid with Parratt and kinematic calculations]{Measured
  near-critical Bi$_2$Se$_3$ \(m=0\) profile overlaid with the two limiting
  calculations used by the optical-to-kinematic handoff. The solid black curve
  is the background-subtracted detector-derived data, the dotted curve is the
  Parratt reflectivity calculation, and the dashed curve is the finite-stack
  kinematic calculation. All three are shown on the same normalized-intensity
  axis over \(0.2\leq L\leq 1.05\), before the \(003\) peak. The common overlay
  makes the low-\(L\) optical regime and the convergence of the calculated
  branches toward the measured high-\(L\) tail directly visible. The smoothly
  blended curve used internally by the model is omitted.}
  \label{fig:parratt_handoff}
\end{figure}
